% ID: SEL-D0010-BC-E05
% ESTADO: validada. S011; tres pasadas del mismo autor, sin segundo revisor.
\textbf{Datos y árbol de probabilidades.}\par\nopagebreak[4]
Sean $E$ el suceso estar enfermo y $+$ el resultado positivo. Los datos son $P(E)=0.15$, $P(+\mid E)=0.92$ y $P(+\mid E^c)=0.04$. Por complemento, $P(E^c)=0.85$.
Las cuatro hojas se calculan multiplicando la probabilidad del primer suceso por la condicionada:
\[
\begin{array}{c|rr|r}
&+&- &\text{Total}\\\hline
E&0.138&0.012&0.15\\
E^c&0.034&0.816&0.85\\\hline
\text{Total}&0.172&0.828&1
\end{array}
\]
\textbf{a) Enfermedad, condicionado al positivo.}\par\nopagebreak[4]
Aplicamos Bayes, con denominador dado por la probabilidad total de positivo:
\[
\boxed{P(E\mid+)=\frac{0.15\cdot0.92}{0.172}
=\frac{69}{86}\approx0.802326.}
\]
No se confunde con $P(+\mid E)=0.92$: son condiciones distintas.

\textbf{b) Enfermedad y resultado negativo.}\par\nopagebreak[4]
Se pide una intersección:
\[
\boxed{P(E\cap-)=P(E)P(-\mid E)=0.15\cdot0.08=0.012.}
\]
\textbf{c) Enfermedad, condicionado al negativo.}\par\nopagebreak[4]
\[
\boxed{P(E\mid-)=\frac{P(E\cap-)}{P(-)}
=\frac{0.012}{0.828}=\frac1{69}\approx0.014493.}
\]
Los denominadores son positivos y las cuatro hojas suman uno. Se han usado únicamente los datos probabilísticos del enunciado.
